% Section 2.1, from lectures/L02/appendix/a_shell_and_processes.md.

\section{The Shell, the Process Tree, and Where Things Live}
\label{c2:app:a}

This section is the first half of this chapter's material: what the shell actually is, what
it asks the kernel to do when you type a pipeline, how processes relate to one another, and what the
filesystem hierarchy is for. The second half, the three pseudo-filesystems through which all of this
is visible, is \S\ref{c2:app:b}.

The framing for the whole chapter is worth stating once. \textbf{Every command you type is a
question you are asking the kernel, or an instruction you are giving it.} The shell is not part of
the kernel and has no special powers; it is an ordinary program that makes the same system calls
your own programs make. Once that is genuinely believed rather than merely accepted, most of the
command line stops needing to be memorised.

\subsection{The shell is an ordinary program}
\label{c2:app:a1}

\file{/bin/sh} on the target is a symlink to BusyBox, which is a normal ELF binary with no
privileges of any kind. It reads a line, works out what you meant, and makes system calls. That is
all.

The consequence people find surprising is that \textbf{most of what a shell appears to do is not
done by the shell at all}. Wildcards are expanded by the shell before the program ever runs, so
\code{rm *} hands \code{rm} a list of filenames and \code{rm} never sees a \code{*}. Redirection is
done by the shell before \code{exec}, so a program has no idea its output is going to a file. And
the program you ran received none of the syntax: by the time \code{main} starts, \code{argv} holds
the expanded, redirection-free result.

There is a short list of things the shell \emph{must} do itself, and the list is short for one
reason: they change the shell's own state, and a child process cannot change its parent. \code{cd}
is the canonical example. If \code{cd} were a program in \file{/bin}, the shell would fork, the
child would change \emph{its} directory, and the child would exit. Nothing would have happened. So
\code{cd} is built in, and so are \code{export}, \code{exit}, \code{umask} and a handful of others.

\begin{console}
$ type cd
cd is a shell builtin
\end{console}

\subsection{What a pipeline actually is}
\label{c2:app:a2}

\code{ps | grep sh | wc -l} looks like one thing. It is three processes, two pipes, and a specific
sequence of system calls. This is the single most useful thing in the chapter to have drawn.

A \textbf{pipe} is a kernel object with two file descriptors: one you write to, one you read from.
\code{pipe()} creates it. It has no name and no presence in the filesystem.

A process gets a new program by \textbf{forking} and then \textbf{exec'ing}: \code{fork()} produces
a copy of the current process, and \code{exec()} replaces that copy's program while keeping its file
descriptors. That last clause is what makes everything work. File descriptors survive \code{exec},
so the shell can arrange a child's descriptors \emph{before} the child becomes the program it will
run, and the program starts life already plumbed.

For the two-stage pipeline \code{A | B}:

\begin{console}
  shell
    |
    +-- pipe()             creates fds [r, w]
    |
    +-- fork() ------------> child A
    |                          dup2(w, 1)      stdout now writes into the pipe
    |                          close(r); close(w)
    |                          exec("A")       A starts, unaware of any of this
    |
    +-- fork() ------------> child B
    |                          dup2(r, 0)      stdin now reads from the pipe
    |                          close(r); close(w)
    |                          exec("B")
    |
    +-- close(r); close(w)   the shell must close its own copies
    +-- wait() for both
\end{console}

\noindent \textbf{The closes are not tidiness.} A pipe's read end returns end-of-file only when
\emph{every} copy of the write end is closed. If the shell keeps its copy open, \code{B} waits
forever for an EOF that never comes, and the pipeline hangs. This is the most common bug in
hand-written pipeline code and it is worth knowing before you write any.

Two consequences follow that explain things people find odd:
\begin{itemize}
  \item \textbf{The stages run concurrently, not in sequence.} \code{A} is not finished when
    \code{B} starts. A pipeline reading a huge file starts producing output immediately.
  \item \textbf{The pipe has a fixed-size buffer}, 64 KB by default. When it fills, the writer
    blocks. That is the entire flow-control mechanism, and it is why \code{yes | head -1} never fills
    memory. What ends it is the reader leaving: once \code{head} exits, nothing holds the read end
    open, and \code{yes}'s next write kills it with \code{SIGPIPE}.
\end{itemize}

\subsection{Exit status, and how the shell knows what happened}
\label{c2:app:a3}

Every process returns one byte to its parent. Zero means success and anything else does not;
\code{$?} holds the last one.

\begin{console}
$ true; echo $?
0
$ false; echo $?
1
\end{console}

\noindent That single byte is the entire contract between a program and a script, and it is why
\code{command && next} and \code{command || fallback} work. In a pipeline, \code{$?} is the status
of the \emph{last} stage only, which means \code{grep foo file | wc -l} reports success even when
\code{grep} found nothing. This bites people writing CI scripts, and the fix is \code{set -o
pipefail}, which this repository's own scripts use for exactly that reason.

\textbf{A signal is reported differently.} A process killed by signal N reports as \code{128 + N},
so a program killed by \code{SIGKILL} (9) gives 137. That number appears in \file{ci/test.sh}, which
uses it to tell a target that hung from a target whose tests failed.

\subsection{The filesystem hierarchy, honestly}
\label{c2:app:a4}

The layout is conventional rather than enforced; the kernel cares about almost none of it.

\begin{booktable}{@{}l>{\hsize=0.85\hsize}L>{\hsize=1.15\hsize}L@{}}
  \thead{Path} & \thead{Holds} & \thead{Notes} \\
  \midrule
  \code{/bin} & Essential user commands & On most systems now a symlink to \code{/usr/bin} \\
  \code{/sbin} & Essential system commands & Same \\
  \code{/lib} & Shared libraries and kernel modules & \code{/lib/modules/$(uname -r)/} \\
  \code{/etc} & Configuration, text, machine-local & Never binaries \\
  \code{/dev} & Device nodes & Populated by devtmpfs, see \S\ref{c2:app:b5} \\
  \code{/proc} & Process and kernel information & Not files, see \S\ref{c2:app:b1} \\
  \code{/sys} & The driver model & Not files either, see \S\ref{c2:app:b4} \\
  \code{/tmp} & Scratch, usually cleared at boot &  \\
  \code{/var} & State that changes: logs, spools & The part of the tree you cannot mount read-only \\
  \code{/usr} & Everything not needed to boot & Historically a separate partition \\
  \code{/root} & Root's home directory & Not \code{/} \\
\end{booktable}

\noindent The \file{/bin} against \file{/usr/bin} split is worth ninety seconds because people
assume it encodes a principle. It does not. In 1971 the Unix authors filled a disk, mounted a second
one at \file{/usr}, and moved the overflow there; the split has been rationalised ever since. Modern
distributions have merged them back, and \file{/bin} is a symlink.

\textbf{What matters for embedded} is which directories must exist before \code{init} runs, and the
answer is short: mount points for whatever you intend to mount, and the path named by
\code{rdinit=}. Our own root filesystem creates \file{/proc}, \file{/sys}, \file{/dev}, \file{/tmp},
\file{/lab} and \file{/root} as empty directories in \file{ci/rootfs.sh}, and that is genuinely all
that is needed. An empty directory is required because \textbf{mounting onto a path that does not
exist fails}, and \file{/init} mounting \file{/proc} is one of the first things that happens.

\subsection{Processes: the tree, and the states}
\label{c2:app:a5}

Every process except PID 1 has a parent, so processes form a tree.

\begin{console}
$ ps -o pid,ppid,stat,comm
  PID  PPID STAT COMMAND
    1     0 S    init
   56     1 S    run.sh
   70    56 R    ps
\end{console}

\noindent \code{STAT} is the state, and four of them matter:

\begin{booktable}{@{}lLl@{}}
  \thead{State} & \thead{Means} & \thead{Consumes CPU} \\
  \midrule
  \code{R} & Running, or ready to run and waiting for a CPU & Yes \\
  \code{S} & Interruptible sleep: waiting for something, killable & No \\
  \code{D} & Uninterruptible sleep: usually waiting on I/O & No \\
  \code{Z} & Zombie: exited, but its parent has not collected it & No \\
\end{booktable}

\noindent \textbf{\code{D} is the one that alarms people.} A process in \code{D} cannot be killed,
not even with signal 9, because it is inside a system call that is not at a point where a signal can
be delivered. It is almost always waiting for hardware. If it is stuck there permanently, something
below it is broken, and the process is a symptom rather than the problem.

\textbf{A zombie is not a leak of anything except a process ID.} When a process exits, its memory is
freed immediately, but its exit status has to be kept until its parent asks for it with
\code{wait()}. That retained status is the zombie. A parent that never calls \code{wait()}
accumulates them, and the fix is always in the parent.

\textbf{An orphan is different and is not a problem.} If a parent exits first, its children are
re-parented to PID 1, which is expected to \code{wait()} for them. That is the mechanism by which
daemons end up owned by init, and it is why \code{PPID} of \code{1} is normal rather than
suspicious.

\subsection{Signals}
\label{c2:app:a6}

A signal is a number delivered to a process, which either has a handler for it, or gets the default
behaviour.

\begin{booktable}{@{}lr>{\hsize=0.52\hsize}L>{\hsize=1.48\hsize}L@{}}
  \thead{Signal} & \thead{Number} & \thead{Default} & \thead{Catchable} \\
  \midrule
  \code{SIGINT} & 2 & Terminate & Yes. This is Ctrl-C \\
  \code{SIGKILL} & 9 & Terminate & \textbf{No} \\
  \code{SIGSEGV} & 11 & Terminate and dump & Yes \\
  \code{SIGTERM} & 15 & Terminate & Yes. The polite one, and what \code{kill} sends by default \\
  \code{SIGSTOP} & 19 & Stop & \textbf{No} \\
  \code{SIGCHLD} & 17 & Ignore & Yes. Sent to a parent when a child exits \\
\end{booktable}

\noindent \code{SIGKILL} and \code{SIGSTOP} cannot be caught, blocked or ignored, which is what
makes them the last resort. Everything else can be handled, and a program that handles
\code{SIGTERM} to flush its state before exiting is a program that behaves correctly when the system
shuts down.

The relevance to the rest of this course is that \textbf{a driver's blocking \code{read()} has to
cope with a signal arriving while it waits.} That is Chapter~\ref{c9:ch}'s subject, and the reason a
user pressing Ctrl-C expects something to happen is this table.

\subsection{PID 1, and two ways it is not an ordinary process}
\label{c2:app:a7}

PID 1 is the first userspace program the kernel starts, named by \code{rdinit=} or \code{init=} on
the kernel command line, or found at one of a few default paths. On our target it is
\file{tools/target/init}, which is a shell script.

It is special in exactly two ways, and both are worth demonstrating.

\textbf{It is immune to signals it has not handled.} The kernel discards any signal sent to PID 1
for which PID 1 has installed no handler, including \code{SIGKILL}. So:

\begin{console}
# kill -9 1
# echo "still here"
still here
\end{console}

\noindent Nothing happens. No error, no message, no dead system. This surprises nearly everyone, and
it is the reason ``just kill init'' is not how you shut a machine down. The rule exists because the
kernel cannot function without PID 1, so it refuses to let anything remove it by accident.

\textbf{It may not exit.} The one thing PID 1 must never do is return. If it does, the kernel stops
immediately:

\begin{console}
init: I am PID 1, and I am about to exit
[    3.183995] Kernel panic - not syncing: Attempted to kill init! exitcode=0x00000000
\end{console}

\noindent That is what a real init's infinite loop is for, and it is why \file{tools/target/init}
ends in either \code{exec /bin/sh} or \code{poweroff -f} and never simply falls off the end.
\textbf{A script that runs to completion as PID 1 panics the machine}, which is a memorable way to
discover you forgot the last line.

The pair together is the point. You cannot kill init, and init must not leave.

\subsection{BusyBox, measured}
\label{c2:app:a8}

The entire userland of this target is one binary. Not one binary per command; one binary:

\begin{console}
# ls -l /bin/busybox
-rwxr-xr-x    1 root     root       2108608 ... /bin/busybox
# busybox --list | wc -l
402
\end{console}

\noindent \textbf{402 commands in 2,108,608 bytes}, statically linked, with no libc alongside it.
Every command in \file{/bin} is a symlink to that one file, and BusyBox decides what to do by
looking at \code{argv[0]}.

\begin{console}
# ls -l /bin/ls
lrwxrwxrwx    1 root     root            12 ... /bin/ls -> /bin/busybox
\end{console}

\noindent Compare against a desktop: GNU coreutils alone is around 5 MB for roughly 100 commands,
before libc, and a full \file{/usr/bin} is hundreds of megabytes. The ratio is not subtle.

\textbf{What it costs is real and specific, and comes in three kinds.}

\emph{Missing applets.} Some tools simply are not there. On this target:

\begin{booktable}{@{}llL@{}}
  \thead{Tool} & \thead{Present} & \thead{Why} \\
  \midrule
  \code{dmesg} & Yes & BusyBox applet \\
  \code{ps}, \code{top} & Yes & BusyBox applets, with fewer options than the ones you know \\
  \code{lsof} & Yes & BusyBox applet, and much reduced: it lists fds and little else \\
  \code{free} & Yes & BusyBox applet \\
  \code{strace} & \textbf{No} & Not a BusyBox applet at all; it is a separate program \\
  \code{ss} & \textbf{No} & Same \\
  \code{vmstat} & \textbf{No} & Same \\
  \code{journalctl} & \textbf{No} & It is systemd's, and this system has no systemd \\
\end{booktable}

\noindent Check for yourself with \code{busybox --list | grep -x vmstat}, which prints nothing.

The last two rows are two different absences and it is worth keeping them apart. \code{strace},
\code{ss} and \code{vmstat} are absent because nobody wrote a BusyBox applet for them and you could
cross-compile and install one. \code{journalctl} is absent because it is part of a different init
system; there is nothing to install, because there is no journal.

\emph{Reduced options.} \code{busybox find} has no \code{-printf}. \code{busybox ps} does not take
most of the flags you know. The applet is there and the flag is not, which is the more annoying
failure because it shows up halfway through a script.

\emph{Different behaviour.} The worst kind, because it is silent. This course has already been
bitten by one: \textbf{BusyBox's \code{rmmod} returns success even when the kernel refused to unload
the module}, which is why Chapter~\ref{c4:ch}'s lab checks whether the module is still loaded
instead of checking an exit status. A test written against coreutils behaviour passes on your laptop
and lies on the target.

The engineering summary is that BusyBox is what you ship and coreutils is what you develop against,
and the difference between them is a category of bug that only appears on the target. Test scripts
on the target, not on your laptop.
